% =====================================================================
%  INTRODUÇÃO / 1 Contextualização
% =====================================================================
\setdiv{INTRODUÇÃO}
\chapter{Contextualização}

Nos últimos anos, a economia brasileira tem enfrentado um cenário desafiador,
marcado por instabilidade macroeconômica e oscilações no mercado financeiro. Para
tentar controlar a inflação pressionada pela alta do dólar e pelos desafios nas
contas públicas, o Banco Central tem buscado contê-la por meio do aumento da taxa
Selic, que alcançou 14,25\% ao ano e pode ultrapassar 15\%.

Definida periodicamente pelo Comitê de Política Monetária (Copom), essa taxa
reflete as condições macroeconômicas do país e as expectativas para a inflação em
2025, segundo previsões de instituições financeiras como Citi, Itaú e XP
\cite{ayres2025}.

Esse aumento da Selic reflete preocupações com a inflação, que vem subindo não só
pelo custo maior de produtos importados, consequência direta da desvalorização do
Real, mas também por fatores internos, como o consumo elevado e a falta de
confiança na estabilidade fiscal do país. Essa instabilidade do mercado financeiro
está diretamente relacionada à política monetária adotada pelo Banco Central, com
destaque para a taxa Selic. Como principal instrumento para o controle da inflação
e a regulação da atividade econômica, a Selic influencia diretamente outras taxas
de juros, como as de empréstimos, financiamentos e aplicações financeiras
\cite{bcbSelic}. Mudanças significativas na Selic podem alterar a dinâmica do
mercado financeiro, impactando tanto o custo do crédito quanto o retorno dos
investimentos.

\begin{table}[htb]
\centering
\caption{Variações da SELIC anual, Meta e Real (2020-2025*)}
\label{tab:selic}
\fontsize{10}{12}\selectfont
\begin{tabular}{ccc}
\toprule
\textbf{Ano} & \textbf{SELIC Meta (\%)} & \textbf{SELIC Real (\%)} \\
\midrule
2025 & 14,25* & -- \\
2024 & 12,25 & 12,15 \\
2023 & 11,75 & 11,65 \\
2022 & 13,75 & 13,65 \\
2021 & 9,25 & 9,15 \\
2020 & 2,00 & 1,90 \\
\bottomrule
\end{tabular}

\vspace{2pt}
{\fontsize{10}{12}\selectfont Fonte: Banco Central do Brasil (2025).}
\end{table}

Apesar do cenário desafiador, a economia brasileira continua crescendo, com um
avanço de 3,3\% em 2024 e previsão de 2,5\% para 2025 no PIB \cite{ayres2024}. No
entanto, esse crescimento acontece em meio a desafios. A alta dos juros torna o
crédito mais caro e dificulta novos investimentos, enquanto a queda do Real, apesar
de favorecer exportadores, encarece produtos e reduz o poder de compra da
população.

Por outro lado, a crescente digitalização e modernização do mercado financeiro
impulsiona o desenvolvimento de soluções tecnológicas que visam otimizar o processo
de tomada de decisão e mitigar os riscos inerentes ao mercado, como o surgimento do
trading algorítmico\footnote{Negociação de ativos automatizada por algoritmos.}
\cite{almeida2016,silva2022}. Nesse contexto, essa ferramenta, que se configura como
uma forma automatizada de investir por meio de algoritmos, ganha relevância na
execução de operações com modelos matemáticos e estatísticos. Essa abordagem busca
proporcionar maior eficiência operacional ao identificar padrões de tendência nos
ativos financeiros. Embora o estudo de \citeonline{almeida2016} aponte desafios em
relação à lucratividade em períodos de crise, os resultados de ambos os trabalhos
demonstram o potencial do trading algorítmico para otimizar a tomada de decisões e
buscar resultados consistentes, sobretudo em cenários de mercado mais estáveis.

Conforme apontado em artigo da Forbes \cite{cojocaru2024}, essa tendência digital,
que já automatiza entre 60\% e 73\% das negociações de ações nos EUA, está se
expandindo globalmente. Empresas como Trade Ideas, TrendSpider e Imperative
Execution utilizam inteligência artificial para analisar grandes volumes de dados,
prevendo tendências de mercado e padrões, e plataformas como TD Ameritrade,
Interactive Brokers e TradeStation oferecem ferramentas e APIs para que
investidores, inclusive de varejo, desenvolvam e implementem estratégias
algorítmicas.

Diante desse cenário, este trabalho desenvolveu uma ferramenta baseada em trading
algorítmico e modelos estatísticos para analisar padrões de tendência no mercado de
ativos, explorando as oportunidades geradas pela volatilidade.

\section{Problema}

A alta inflação é constante no mercado brasileiro da atualidade, amplamente estudada
e divulgada pelo BCB, mais especificamente pelo órgão COPOM, que regula e configura
a taxa mais adequada para os níveis inflacionários no país. Os sucessivos aumentos
da taxa básica de juros tornam investimentos menos favoráveis e comprometem o poder
de compra, representando um desafio central para a manutenção de patrimônio.

A instabilidade inerente aos ativos financeiros dificulta a realização de previsões
confiáveis, enquanto a persistente inflação desvaloriza a moeda nacional, exigindo
dos investidores a adoção de estratégias que visem primordialmente a proteção do
capital e, secundariamente, a obtenção de retornos reais positivos.

Diante desse cenário macroeconômico adverso, torna-se essencial a busca e adoção de
modalidades de investimento que sejam não apenas adaptadas à inflação vigente, mas
também resilientes às frequentes e, por vezes, abruptas oscilações do mercado. A
complexidade observada nos modelos tradicionais de análise de investimento, quando
confrontados com a dinâmica atual do mercado, amplia o desafio. Emerge, portanto, a
necessidade premente do desenvolvimento de ferramentas analíticas avançadas que
possam auxiliar na identificação de padrões de tendência e na tomada de decisões de
investimento mais informadas e potencialmente mais seguras neste contexto econômico
particular.

\section{Hipóteses}

\subsection{Hipótese Central}
Este trabalho partiu da premissa de que a integração metodológica de previsão de séries temporais com modelo
ARIMA e modelagem de volatilidade condicional com modelo GARCH em uma ferramenta algorítmica unificada, desenvolvida em MQL5, permite a geração
de sinais de negociação para ativos com maior eficácia, aferida em termos de
potencial de acerto e adequação ao risco, do que abordagens que empregam estas
técnicas de forma isolada ou incompleta. A viabilidade de combinar diferentes
ferramentas de análise para obter resultados superiores encontra respaldo em
trabalhos como o de \citeonline{batista2024}, que demonstrou resultados positivos ao
integrar indicadores técnicos em um modelo de negociação.

\subsection{Hipótese da Complementaridade Preditiva e de Risco}
Partiu-se da premissa que a combinação sequencial dos modelos ARIMA e GARCH constitui um avanço
significativo em relação a análises baseadas em Médias Móveis, como as exploradas
por \citeonline{almeida2016}. Especificamente, esperava-se que: o modelo ARIMA capturasse
a dinâmica temporal e fornecesse previsões de curto prazo da direção dos preços com
maior refinamento do que as Médias Móveis. A subsequente aplicação do modelo GARCH
aos resíduos gerados pelo ARIMA incorpora com maior detalhe a modelagem da
volatilidade e do risco associado às previsões, permitindo uma avaliação mais ampla
do cenário de mercado.

\subsection{Hipótese da Adequação ao Risco}
Partiu-se da premissa que a modelagem estatística da volatilidade via GARCH, complementada por métricas quantitativas de risco, assegura uma avaliação mais robusta das condições de mercado. A análise integrada de tendência (ARIMA) e risco (GARCH) reduz a probabilidade de decisões inadequadas em momentos de mercado desfavorável, aumentando a seletividade e a segurança das operações identificadas pela ferramenta.

\section{Objetivo}

\subsection{Objetivo Geral}
Desenvolver e avaliar uma ferramenta algorítmica integrada, utilizando a linguagem de programação
MQL5 e a plataforma MetaTrader 5, para análise e previsão de tendências no mercado
de ativos, com o propósito de gerar retornos superiores à taxa Selic, com foco na economia brasileira.

\subsection{Objetivos Específicos}
\begin{itemize}
  \item Aprofundar conceitos teóricos do mercado financeiro, com ampla ênfase em
        volatilidade e tendências;
  \item Compreender e aplicar os conhecimentos na plataforma MetaTrader em junção
        com MQL5 na automação de estratégias;
  \item Implementar modelos estatísticos (ARIMA e GARCH) para previsão de séries
        temporais e modelagem de volatilidade;
  \item Comparar resultados obtidos com a taxa Selic vigente;
  \item Documentar obstáculos encontrados e lições aprendidas durante a execução do projeto.
\end{itemize}

\section{Justificativa}

Esta pesquisa se justifica inicialmente pela revisitação e aplicação contemporânea
do conceito explorado por \citeonline{almeida2016}. Seu estudo demonstrou a
viabilidade do uso de robôs de investimento via MetaTrader em diferentes períodos
sobre o ativo PETR4. Contudo, houve uma evolução significativa nos cenários
econômicos e nas tecnologias disponíveis desde então, o que motiva esta nova
abordagem, alinhada às suas próprias recomendações futuras para aplicação em
mercados de alta volatilidade.

Ademais, por consequência direta do presente cenário econômico de inflação
persistente e juros elevados no Brasil, outro ponto central que justifica este
trabalho é a necessidade de identificar e avaliar alternativas de investimento que
possam se mostrar viáveis e potencialmente lucrativas. Este estudo visa contribuir
para mitigar os efeitos do congelamento financeiro que pode afetar investidores
nesse contexto, oferecendo análises que possam embasar decisões para aqueles
interessados em diversificar seus portfólios.

Além disso, a pesquisa busca agregar valor ao conhecimento acadêmico na área de
finanças quantitativas e trading algorítmico no contexto brasileiro, podendo servir
de base ou inspiração para futuros projetos relacionados.

\section{Natureza e Tipo da pesquisa a ser desenvolvido}

A presente pesquisa se enquadra, inicialmente, na classificação de ciência empírica,
uma vez que se dedica à investigação de fenômenos observáveis do mundo real.
Conforme \citeonline{wazlawick2014}, a ciência empírica se caracteriza pela busca de
conhecimento através da experiência e da observação, o que se manifesta neste estudo
na análise de dados históricos de preços de ações para a identificação de padrões e
tendências. A coleta e o tratamento desses dados são fundamentais para a construção
e validação da ferramenta algorítmica.

Adicionalmente, a pesquisa possui um caráter aplicado. \citeonline{wazlawick2014}
destaca que a ciência aplicada se preocupa com a utilização prática do conhecimento
científico para a solução de problemas específicos. Neste contexto, o
desenvolvimento de uma ferramenta para auxiliar investidores no mercado de ativos
demonstra claramente a intenção de aplicar os conhecimentos gerados para atender a
uma demanda concreta, buscando otimizar o processo de tomada de decisão e a gestão
de riscos.

No que tange à sua natureza, a pesquisa se aproxima de uma ciência exata, embora com
ressalvas. A utilização de modelos estatísticos como ARIMA e GARCH, bem como a
adoção de métricas quantitativas para avaliar o desempenho da ferramenta, refletem a
busca por precisão e objetividade, características marcantes das ciências exatas. No
entanto, é fundamental reconhecer que o mercado financeiro, objeto de estudo desta
pesquisa, é intrinsecamente marcado pela incerteza e pela volatilidade. Essa
característica introduz um elemento de inexatidão que diferencia este campo de
investigação de sistemas puramente determinísticos, nos quais as ciências exatas
tradicionalmente se concentram.

Sob a perspectiva do rigor metodológico, a pesquisa busca alinhamento com uma ciência
hard. Segundo \citeonline{wazlawick2014}, as ciências hard se caracterizam pela
utilização de métodos formais e quantitativos, pela busca de leis gerais e pela
ênfase na objetividade e na replicabilidade. A aplicação de métodos estatísticos, a
definição de métricas de avaliação e a validação empírica por meio de backtesting
demonstram a preocupação em conferir solidez e validade aos resultados obtidos.

Por fim, a pesquisa possui um componente nomotético. \citeonline{wazlawick2014}
explica que as ciências nomotéticas buscam estabelecer leis gerais ou princípios que
expliquem classes de fenômenos. Neste estudo, a análise de tendências e volatilidade
no mercado de ações tem como objetivo identificar padrões que possam, dentro de
certas limitações, auxiliar na previsão do comportamento futuro dos ativos. Embora a
previsão no mercado financeiro seja inerentemente complexa, a busca por
regularidades e a tentativa de formular generalizações demonstram o caráter
nomotético da investigação.

Em relação aos objetivos, a pesquisa se classifica como explicativa. Conforme
\citeonline{wazlawick2014}, a pesquisa explicativa busca identificar as causas ou
razões de determinados fenômenos, aprofundando o conhecimento sobre as relações
entre variáveis. Neste caso, o estudo buscou explicar como a combinação de diferentes
modelos estatísticos pode influenciar a capacidade preditiva e a gestão de riscos no mercado de ações, investigando os mecanismos pelos
quais essa integração otimiza o desempenho da ferramenta algorítmica.

No que concerne aos procedimentos técnicos, a pesquisa se desenvolveu por meio de
pesquisa bibliográfica e pesquisa experimental. A pesquisa bibliográfica foi
fundamental para a revisão da literatura existente sobre modelos estatísticos
e trading algorítmico, fornecendo a base teórica e conceitual
para o estudo. A pesquisa experimental, por sua vez, foi realizada por meio de
backtesting na plataforma MetaTrader 5, envolvendo a manipulação de variáveis
(parâmetros dos modelos) e a observação dos resultados (taxa
de acerto, lucro/prejuízo) para validar a eficácia da ferramenta desenvolvida.

\section{Organização dos capítulos}

O Capítulo 1 apresenta a Fundamentação Teórica, oferecendo a base de conhecimento
necessária para a compreensão do projeto.

Inicialmente, foram abordados conceitos introdutórios que contextualizam a pesquisa.
Isso inclui partes sobre economia e princípios básicos do mercado
financeiro, como o funcionamento de corretoras, bolsas de valores e as
características dos ativos negociados.

Mais adiante, o capítulo se aprofunda em conceitos mais técnicos, importantes para a
análise e modelagem. Foram discutidas as diferentes abordagens da análise técnica de
ativos, com foco nos métodos estatísticos. A plataforma MetaTrader 5 e sua linguagem
MQL5, utilizadas no desenvolvimento, foram apresentadas, assim como os tipos de
indicadores de mercado e os modelos estatísticos para análise de séries temporais e
volatilidade.

Por fim, o capítulo apresenta uma revisão de pesquisas anteriores. Esta seção
relaciona o presente estudo com trabalhos já existentes na área,
identificando como esta pesquisa se posiciona e contribui para o campo.

O Capítulo 2 detalha a Metodologia utilizada, explicando o processo de construção e
validação da ferramenta algorítmica.

Este capítulo começa descrevendo a natureza da pesquisa e o ambiente para os
experimentos. Foram especificadas as ferramentas de software,
linguagens de programação e as fontes de dados para o ativo financeiro em estudo,
incluindo como os dados do mercado brasileiro foram acessados.

Em seguida, é descrito como o experimento de pesquisa foi conduzido, com a descrição
das etapas para a implementação dos diferentes componentes da ferramenta. O processo
iniciou com as análises baseadas em indicadores de tendência, avançando para a modelagem
estatística dos preços e a análise da volatilidade dos resíduos via ARIMA e GARCH.

O capítulo apresenta também os obstáculos encontrados durante a execução, conclui
estabelecendo os critérios para a avaliação dos modelos desenvolvidos e da ferramenta
como um todo, detalhando as métricas de desempenho utilizadas nos testes.
Adicionalmente, são apresentados os resultados preliminares obtidos e as limitações
bem como perspectivas para trabalhos futuros.
